\scanpage{31}\mspage{30}
% @ top, part a
a) $F'$ est une figure. \deleted{i.e.} Notons que \deleted{$|F'|=$}
% @ top, support with braces
\[
 |F'|=\underbrace{|G'|}_{|G|}\cup\underbrace{|F\setminus G|}_{\supset |F|\setminus |G|}\deleted{\unsure{=|\supset}}\supset |F|
\]
% @ upper part, connector
donc, comme $|F'|\subset |F|$, on a
% @ upper part, support relation: hooked upper stroke over detached lower bar
\[
 |F'|\unsure{\supseteq|=}|F|
\]
% @ upper part, argument
Il faut prouver que $\forall x\in |F'|=|F|$, $\exists$ plus petit \unsure{élémt|élément|élt} de $F'$ qui le contienne. Distinguons deux cas

% @ upper middle, first case
1$^\circ$) $x\in |G|=|G'|$, soit $A\in G'\deleted{\unsure{\circ}}$ tel que $x\in A^{\circ(G')\deleted{\unsure{)|?}}}$. Si $B\in F'$ et \deleted{$B\ni$} $x\in B$, prouvons $A\subset B$. Si $B\in G'$, c'est évident. Si $B\in F\setminus G$, alors \deleted{$B$ contient} $A^{\circ(G')}\cap B\ne\varnothing$ donc $A^{\unsure{\circ|\bullet}}\subset B$ OK

% @ middle, second case
2$^\circ$) $x\in |\deleted{\unsure{G}}F|\setminus |G|$, alors les $\deleted{B}A\in F'$ qui contiennent $x$ ne sont pas dans $G'$, ils sont donc dans $F\setminus G$. D'autre part, l'unique \unsure{élémt|élément|élt} $A$ de $F$ tel que $x\in A^{\circ(F)}$ n'est pas dans $G$ puisque $x\notin |G|$, il est donc dans $F\setminus G$, et \deleted{\unsure{si $x\in A$}} c'est lui plus petit él. de $F$ (et à fortiori de $F\setminus G$) qui contient $x$, OK~!

% @ lower middle, part b
b) $F'$ \deleted{raffine}\inserted{subdivise} $F$. Comme $|F'|\unsure{\supseteq|=}|F|$, il faut voir seult que $F'$ raffine $F$. Donc \deleted{que \unsure{tout}} que
% @ lower part, cancelled attempt
\deleted{1$^\circ$) \deleted{$\forall$} $A\in F'$, 
% @ upper interlinear addition within the cancelled draft
\inserted{l'ens $F^A$ des $B\in F$ tels que $A\subset B$}
% @ base line: a single short word precedes contenu, not two separate et/est occurrences
\unsure{est|et} \unsure{contenu dans} un plus petit
% @ left slanted addition, with a leader toward the following abbreviation
\inserted{a un plus petit} \unsure{élémt|élément|élt} de $F$. C'est évident si $A\in F$ c'est évident. Si}
\deleted{\unsure{str.} de de $F'$ raffine $F$. Si}

% @ lower part, restated first condition
1$^\circ$) \deleted{$\forall$} $A\in F'$, l'ens $F^A=\{B\in F\mid \deleted{B}A\subset B\}$ a un plus petit \unsure{élémt|élément|élt}. C'est évident si $A\in F$ \deleted{$(B=A)$} (on prend $B=A$) et si $A\in G'$, c'est vrai puisque $G'\ll F$.

% @ bottom, second condition
2$^\circ$) Si $A\in F'$, $B\in F$, alors $A\cap B$ est \deleted{\unsure{$\mathcal R_{F'}$-satu.|$\mathscr R_{F'}$-satu.}}\inserted{fermé dans $F'$}, i.e. réunion de strates de $F'$. Si $A\in G'$, comme $G'\ll F$, $A\cap B$ est réunion de strates de $G'$, qui sont aussi des strates de $F'$.
